\section{考察および今後の展望}
\label{sec:discussion_future}

本研究では，軟弱地面において足の沈み込みが生じる状況を対象とし，DS から SS への遷移時に発生する転倒挙動の要因を解析するとともに，重心（CoM）高さを制御することによる安定化手法を提案した．シミュレーション結果から，沈み込みが発生する条件において CoM 高さを適切に調整することで，支持脚足首関節に作用するロール方向トルクを低減できること，ならびに SS 期間における ZMP の発散を抑制し，次の一歩の安定性を向上させられることが示された．

これらの結果は，沈み込みによって生じる姿勢の不安定化に対して，CoM 高さが重要な調整パラメータとして機能することを示している．すなわち，床への踏み込みによる沈み込み量を適切に計測・推定することができれば，その情報を基に CoM 高さを必要最小限だけ変化させることで，大規模なモデル拡張や高次の制御器を導入することなく，コンパクトかつ効果的な安定化制御が実現できる可能性がある．この点は，既存の LIPM ベース歩行制御器との親和性が高いという本手法の特徴を強く示すものである．

一方で，本研究の結果から，提案手法には適用可能な条件範囲が存在することも明らかとなった．沈み込み量や沈み込み速度が過度に大きい場合には，CoM 高さ制御が十分に作用する前に姿勢の不安定化が進行し，転倒を回避できない場合がある．このことから，提案手法を実環境へ適用するためには，床の物理特性および沈み込み挙動を考慮した制御指針の整理が不可欠である．

今後の展望としては，沈み込む床の剛性や減衰といった物理特性と，それに伴って必要となる CoM 高さ変化量との関係を定式化することが挙げられる．具体的には，床特性および沈み込み量を入力とし，安定性を維持するために必要な CoM 高さ補正量を出力とする関係式や制御則を導出することで，より汎用的かつ自律的な軟弱地面適応歩行の実現が期待される．また，このような定式化は，実機ロボットへの適用や，異なる床環境への一般化に向けた重要な基盤となると考えられる．
